\begin{abstract}
\addcontentsline{toc}{chapter}{Abstract}
The modern university operates as a high-dimensional, data-intensive ecosystem requiring continuous information retrieval, document indexing, constraint satisfaction, and high-throughput computational analytics. Traditional linear and simple tree-based data structures from introductory curricula are provably insufficient to resolve sub-quadratic text search, sequence alignment, capacity-constrained resource distribution, or computationally intractable scheduling problems.

This report presents the complete design, theoretical analysis, and architectural implementation of \textbf{EduTrack: An Intelligent University Academic Platform}. Engineered to satisfy the rigorous syllabus requirements of Advanced Data Structures and Algorithms (DSA-3), EduTrack encompasses six core advanced algorithmic families:
\begin{enumerate}
    \item \textbf{Exact and Multi-Pattern String Algorithms:} Knuth-Morris-Pratt (KMP), Z-Algorithm, Rabin-Karp polynomial rolling hash with double-hashing, and Aho-Corasick multi-pattern automaton for instant catalog querying and academic taxonomy tagging.
    \item \textbf{High-Dimensional Suffix Structures \& Document Similarity:} Suffix Arrays via prefix doubling, linear-time Suffix Array Induced Sorting (SA-IS), Kasai's linear Longest Common Prefix (LCP) algorithm for cross-assignment plagiarism detection, and Suffix Automata (SAM) for sub-linear substring query evaluation.
    \item \textbf{Advanced Dynamic Programming:} Wagner-Fischer Levenshtein distance, transposition-aware Damerau-Levenshtein distance for typographical query correction, Matrix-Chain Multiplication (MCM) for multi-stage student analytics projection, Bitmask DP for Traveling Academic Auditor venue scheduling ($O(2^n \cdot n^2)$), and Optimal Binary Search Trees (OBST) for search frequency optimization.
    \item \textbf{Network Flow \& Duality:} Residual flow networks, augmenting path methods (Ford-Fulkerson and Edmonds-Karp $O(VE^2)$), Dinic's blocking flow algorithm ($O(V^2E)$), Bipartite Matching for faculty-to-course allocation, and König's Theorem for minimum faculty-course conflict bottleneck analysis.
    \item \textbf{Computational Intractability, Reductions \& Approximation:} Davis-Putnam-Logemann-Loveland (DPLL) Boolean Satisfiability (SAT) solver for examination timetabling constraints, the complete polynomial-time reduction chain ($\text{3-SAT} \le_p \text{CLIQUE} \le_p \text{INDEPENDENT-SET} \le_p \text{VERTEX-COVER}$), and a maximal-matching-based 2-Approximation algorithm for exam invigilator/proctor minimization.
    \item \textbf{Randomized, Streaming \& Parallel Algorithms:} Randomized QuickSort for student merit ranking with expected $O(N \log N)$ complexity, Reservoir Sampling (Algorithm R) for streaming activity log audits, Miller-Rabin primality testing with 64-bit modular arithmetic for cryptographic student token verification, Blelloch's work-efficient Parallel Prefix Scan ($O(N)$ work, $O(\log N)$ span) for cumulative academic credit metrics, multi-threaded Parallel Tree Reduction for batch GPA aggregation, and Brent's Theorem parallel speedup and efficiency modeling.
\end{enumerate}

A foundational requirement governing this system is the \textbf{strict prohibition of \texttt{java.util.*}} inside the algorithm engine core. Every collection, dynamic array, linked list, FIFO queue, LIFO stack, hash table, priority queue, and disjoint set structure is constructed entirely from first principles.

The platform is evaluated against realistic university datasets including 25 multi-disciplinary courses, 20 detailed student records, 12 faculty profiles with workload constraints, multi-paragraph student assignments featuring intentional plagiarism test cases, and a continuous high-frequency activity log stream. Both an interactive numbered Command-Line Interface (CLI) and a comprehensive Java Swing Graphical User Interface (GUI) desktop application are provided. Experimental validation demonstrates that all 16 algorithmic verification benchmarks pass with zero failures, confirming the theoretical correctness, robust design, and practical scalability of the EduTrack architecture.
\end{abstract}
\newpage
